% mage-auto-heading: agreement-is-not-correctness
% mage-auto-heading: example-siemens-other
% mage-auto-heading: the-semantic-gap
% mage-auto-heading: the-earliest-decidable-boundary
% Authoritative LaTeX manuscript; Markdown is a generated compatibility projection.

% <!-- part-title: Alignment -->

% <!-- chapter-title: From Engineering Intent to Enforceable Obligation -->

\chapter{From Engineering Intent to Enforceable Obligation}

% <!-- section-terms: independent-obligation, semantic-gap, agreement-vs-correctness -->

% <!-- point: enforcement-needs-an-independent-obligation-and-a-legible-boundary | Enforcement needs an obligation independent of the work being judged, and a boundary at which that obligation is legible. | terms: independent-obligation, semantic-gap -->

Enforcement requires two things: \textbf{an obligation independent of the work being judged and a boundary where the evidence needed to evaluate that obligation is available.}
The architectural model requires \texttt{Reporting} to obtain billing data through \texttt{BillingAPI}, preserving \texttt{BillingDB} behind that interface.
The job-lifecycle model specifies how events move a job from \texttt{IDLE} through \texttt{RUNNING} to \texttt{DONE} or \texttt{FAILED}.
These are accounts of intended construction and behavior against which an implementation can be evaluated.

% <!-- point: an-obligation-may-live-in-a-model-or-directly-in-a-type-or-policy | An obligation may live in a model, or directly in a type, permission, schema, test, or policy; Alignment does not require a particular representation. | terms: independent-obligation, modeling-principle, alignment-principle -->

The environment needs evidence that corresponds to those meanings.
A dependency analysis can look for paths from \texttt{Reporting} to \texttt{BillingDB} that bypass \texttt{BillingAPI}; a model-based test can submit job events and compare the resulting states with the specified transitions.
The engineer must choose where the relevant evidence becomes available and how the finding should affect the work.
A check may inform a later decision, or a gate may use its result to refuse an action or candidate.
The obligation can also reside directly in a type, permission, schema, test, or policy.
What matters is that it is distinguishable from the implementation being judged and that a mechanism can enforce the resulting decision.

% <!-- point: software-engineering-has-a-traceability-heritage-alignment-builds-on | Requirements traceability connects intent, implementation, and evidence to support evaluation and change analysis. | terms: independent-obligation, traceability -->

Once an obligation is explicit, engineers need to determine whether the implementation satisfies it. You may be familiar with this work from requirements traceability, which connects a requirement to relevant design elements, implementation, tests, and other evidence~\citep{estefan2008mbse}. These connections help engineers examine the design choices made to meet the requirement and find the evidence needed to evaluate the implementation. When a requirement or design changes, the links help identify which parts of the implementation and which checks need to be reconsidered.

% <!-- box-family: inset -->

% <!-- inset-domain: SOFTWARE ENGINEERING -->

% mage-box-prefix: Inset —

% <!-- print-keep-with-next: 20 -->

\begin{bookbox}{inset}{Assurance Cases}
\label{inset-assurance-cases}
In safety-critical engineering, an assurance case makes explicit why a system should be trusted. It connects claims about the system to supporting evidence through an argument explaining why that evidence supports the claims~\citep{iso15026-2-2022,kelly2004gsn}. This goes beyond requirements traceability. Knowing which tests correspond to a requirement does not establish that those tests adequately verify it.

Suppose DocAble claims to preserve document content during accessibility remediation. Comparing the text before and after processing provides some evidence, but is it sufficient? Text extraction might miss content in images, equations, or annotations. An assurance argument would identify these limitations, explain which additional checks address them, and state what remains unverified. Passing every available test does not establish a claim broader than those tests support.

Assurance cases help engineers justify why a system should be trusted. MAGE helps establish and preserve selected properties on which that justification depends.

If you work under a safety standard such as IEC 61508 or DO-178C, MAGE can complement your existing assurance process~\citep{iec61508-1-2010,rtca-do178c-2011}. The claims and arguments in an assurance case can guide what you model, what you check, and what evidence you retain. Agents can help produce and maintain that evidence as the system evolves. The assurance case still needs independent scrutiny. Neither an agent's report nor a passing check establishes that the overall argument is sound.
\end{bookbox}

% <!-- point: traceability-connects-intent-realization-evidence-without-enforcement | Traceability can relate an obligation, its realization, and its evidence without any mechanism enforcing the obligation — a requirement can trace cleanly to a test whose failure is merely reported, and a permission boundary can enforce firmly without a rich traceability model. | terms: independent-obligation, traceability, alignment-principle -->

Traceability can connect an obligation, its realization, and its evidence without any mechanism enforcing the obligation. A requirement may trace cleanly to a test whose failure is merely reported. Conversely, a permission boundary may enforce firmly without a rich traceability model. Alignment therefore separates the relation among intent, realization, and evidence from the decision to enforce the resulting judgment.

% <!-- point: tolerance-independence-comes-from-intent-not-realization | Where an obligation permits variation, the acceptable region must come from governing intent, not from the realization being judged. | terms: independent-obligation, tolerance -->

Where an obligation permits variation, the acceptable region must likewise come from the engineering intent that governs the work. A measured implementation can establish what latency it achieved or what dependencies it contains; whether those values are acceptable requires an independent engineering obligation.

% <!-- point: those-requirements-produce-the-two-lessons-agreement-not-correctness-and-placement-follows-semantics | Two lessons: agreement is not correctness, and placement must follow the meaning of the obligation. | terms: agreement-vs-correctness, semantic-gap -->

A description and the implementation it describes can therefore agree while both violate the intended property. Placement matters as well: if an obligation is enforced at a boundary too small to make it legible, the checker may inspect the wrong thing.

\section{Agreement Is Not Correctness}\label{heading:agreement-is-not-correctness}

% <!-- section-terms: agreement-vs-correctness, independent-obligation -->

% <!-- point: a-drift-check-establishes-agreement-not-that-the-correspondence-is-right | A drift check establishes agreement between two artifacts, not that the correspondence expresses the right intent. | terms: agreement-vs-correctness -->

% <!-- index-def: agreement-vs-correctness -->

% <!-- index-def: mirror-vs-spec -->

% <!-- index-def: drift-caveat -->

A drift check can establish a declared correspondence between an engineering model and its implementation. It cannot establish that the correspondence expresses the right intent. Suppose an endpoint is required to be authenticated. If an independently authored policy says so, a check can compare the implementation against that obligation and flag an open endpoint. If instead the representation is derived only from the implementation, it may faithfully record the endpoint as open. Model and code agree; the security obligation is still violated.

% <!-- point: the-distinction-is-description-versus-independent-obligation | The distinction is not derived versus authored; it is description versus independent obligation. | terms: independent-obligation, agreement-vs-correctness -->

% <!-- index-def: independent-obligation -->

The distinction is between description and independent obligation. A model extracted from code can supply the current \emph{is}; an independently authored invariant can supply the governing \emph{ought}. Authorship alone establishes neither role: a hand-authored model may merely repeat the implementation\textquotesingle s assumptions, while a derived model can faithfully describe its present behavior. Enforcement requires an obligation against which the work can actually be wrong.

% <!-- point: drift-checking-has-a-precise-job-preserve-the-correspondence-the-model-claims | Drift checking has a precise job: preserve a declared correspondence between model and implementation; correctness needs another claim. | terms: agreement-vs-correctness, independent-obligation -->

Drift checking therefore has a precise job: preserve a declared correspondence between model and implementation. If they diverge, surface the disagreement. Conformance asks whether the implementation satisfies an independent engineering obligation, which may be represented in that same model. Acceptance is a separate decision: the Governed Engineering Environment (GEE) determines whether the implementation may proceed, considering applicable controls and evidence beyond those associated with any one model.

% <!-- point: proof-strength-applies-to-the-proposition-proved-not-to-its-adequacy | Proofs and model checks evaluate stated properties under their assumptions; they do not establish that the chosen properties capture the engineering requirement. | terms: agreement-vs-correctness, independent-obligation -->

Stronger verification does not remove this distinction. A proof assistant such as Lean checks proofs that formally stated propositions follow from their assumptions~\citep{lean-reference}. TLA+ lets engineers specify system behavior; the TLC model checker can explore a model to check whether it satisfies stated properties~\citep{lamport2002}. Neither approach establishes that the chosen property captures the engineering requirement. If an agent chooses a property merely because its implementation satisfies it, a valid proof or successful model check may leave the intended obligation untested. Engineers can ask agents to help formalize requirements and develop proofs or checks while retaining responsibility for whether those requirements express what the system must preserve.

% <!-- point: docable-requested-temporal-checking-for-an-obligation-its-models-already-carried | Fix the required behavior before choosing an analysis; independence concerns the source of the obligation and the implementation judged against it. | terms: independent-obligation, agreement-vs-correctness -->

For example, an engineer can require a queueing service to bring every accepted job to a terminal state, then ask an agent to develop an analysis of that requirement. The agent may help formalize the obligation and construct the checker. What matters is that the obligation comes from what the service must do, rather than being chosen merely because the current implementation satisfies it.

% <!-- point: an-obligation-can-test-one-of-three-relations | Three different relations matter — correspondence, conformance, or acceptance — and they are not interchangeable. | terms: independent-obligation, agreement-vs-correctness -->

\autoref{three-relations} shows these three judgments over one implementation. The engineering model is one part of the GEE, whose admission decision may depend on additional obligations and evidence.

\begin{figure}[htbp]
\centering
\includegraphics[width=\linewidth]{book/latex/assets/three-relations-ac66d52458.pdf}
\caption{\emph{Model agreement, obligation satisfaction, and admission.} Correspondence concerns agreement between an engineering model and its implementation; conformance concerns satisfaction of an independent obligation represented in the model. Acceptance is the GEE\textquotesingle s decision to permit the implementation to proceed, based on applicable controls and evidence that may extend beyond the model shown. Evidence for one relation does not establish the others.}
\label{three-relations}
\end{figure}

% <!-- worked-examples: binds-model-to-territory -->

\subsection{Example --- Siemens \{other\}}\label{heading:example-siemens-other}

The Siemens reconstruction describes models, implementation, and engineering analyses kept in correspondence across a digital thread without making every disagreement a build-blocking event. In that account, a human confirms the bill of materials while physics solvers carry parts of the evaluation. The example is useful precisely because correspondence and gating are separate design choices: the coupling can hold without a hard gate deciding the instant the two diverge.

% <!-- worked-examples-end -->

% <!-- point: execution-consumes-the-computation-model-so-it-governs-composition | DocAble's execution machinery consumes a computation model, while an analytical view supports separate analysis and checks; the representation's governing role follows from its consumers. | terms: agreement-vs-correctness, independent-obligation -->

% <!-- box-family: inset -->

% <!-- inset-domain: Case study -->

% mage-box-prefix:

\begin{bookbox}{inset}{Worked inset --- DocAble: A Model That Directs Work and a View That Describes It}
\label{inset-docable-correspondence}
DocAble repairs accessibility problems in Office and PDF documents through a sequence of computations. A DOCX table-header repair, for example, establishes information needed by a later table-structure check. The scheduler must preserve that dependency even though the two operations have separate implementations. An explicit computation model records such relationships. The scheduler reads the model to determine what is ready to run, while each computation's implementation determines how it performs its work.\magefootnote{static-graph}{The computation model describes the pipeline rather than the history of any one remediation. Runtime telemetry can now be joined to computation identities for aggregate measurements such as latency and cost, but the per-session edit record still names realized document mutations rather than computation nodes. DocAble has not yet built a per-run execution graph joining that edit history to the computation model.}

A separate analytical graph makes that composition available for queries and checks. For Office formats, its dependency edges are projected from registry read/write metadata also used by scheduling. For PDF, the analytical relations come from typed declarations of what computations produce and consume, with blocking parity checks maintaining correspondence to the executable declarations. The analytical graph supports inspection and checking, but does not govern execution. The scheduler consumes the computation model, not the analytical graph.

The checks also have a coverage limit. In the examined Office implementation, one live pass still ran outside the registration surface that the parity checks could enumerate. The checks reported agreement for the represented surface, but remained audit-only. Making them blocking would not have closed that gap: the same green result would still have omitted the unregistered pass. The missing correspondence first had to be understood and addressed before the check could justify a broader acceptance claim.

The case illustrates the three judgments in \autoref{three-relations}:
\begin{itemize}
\item \textbf{Correspondence (model agreement).} Parity checks compare the analytical graph with its sources: registry metadata for Office computations and typed declarations for PDF computations. They establish agreement over the represented relationships, not that every executing computation is represented. The unregistered Office pass exposed this limitation.
\item \textbf{Conformance (obligation satisfaction).} The computation model specifies dependencies and composition rules that the scheduler must respect. Whether execution satisfies those obligations is a separate question from whether the analytical graph agrees with its sources. Parity alone does not establish that the scheduler preserves the required computation order.
\item \textbf{Acceptance (permission to proceed).} The GEE can require correspondence or conformance evidence before admitting a change. Making the existing parity checks blocking, however, would not have detected the unregistered pass: it lay outside their coverage. Enforcement cannot compensate for evidence that omits the operation being governed.
\end{itemize}
\end{bookbox}

Traceability can do more than help engineers locate the implementation of a modeled property. When model elements are linked to code and execution records, those relationships can be checked as the system changes. DocAble uses static references to connect model elements to implementation declarations and runtime traces that identify the modeled operations being executed. Static checks can detect when a referenced implementation element is deleted or renamed. Dynamic checks can compare observed operations with the model, revealing modeled operations that were not observed or executions that cannot be associated with a modeled element. These checks complement guidance mechanisms such as dynamic context injection (\magevalue{sec:the-agent-stack}): agents can receive relevant model information while working, and the environment can independently check correspondence as the implementation changes.

The appropriate evidence depends on the model and the property being checked. Architectural models often admit static checks of dependencies, interfaces, and component boundaries. Behavioral models may require execution traces to establish which states or transitions actually occurred. These are not exclusive categories. Static analysis can check declared state-machine structure, while runtime instrumentation can expose architectural interactions that source-level checks miss. In each case, the traceability relationships make the relevant evidence available for comparison with the model.

These checks turn traceability relationships into Alignment mechanisms. A change that breaks a model–implementation relationship can produce a finding without requiring an engineer to rediscover the relationship manually. Static references establish that implementation elements still exist, while runtime traces provide evidence about their execution. Neither establishes that the implementation still behaves as intended, nor that the traceability relationships cover every relevant operation. Where the evidence is adequate, a gate can require the checks to pass before accepting a change.

\section{The Semantic Gap}\label{heading:the-semantic-gap}

% <!-- section-terms: semantic-gap, earliest-legible-boundary, model-drift -->

% <!-- point: an-obligation-can-be-perfectly-stated-and-unenforceable-at-the-boundary-you-check | An obligation can be stated clearly and still be unenforceable at the boundary where you try to check it. | terms: semantic-gap, earliest-legible-boundary -->

A second problem is placement. An obligation can be stated clearly and still be unenforceable at the boundary where you try to check it. The required evidence may span several changes, or the property may emerge only when components operate together. A semantic gap occurs when the chosen boundary does not expose enough information to decide whether the obligation holds. \autoref{semantic-gap} illustrates two such cases. Although both arise during software development, they require checks at different boundaries.

% <!-- sidenote -->

\begin{quote}
\emph{A footnote on where the placement rule comes from.} It echoes the end-to-end argument in system design: a function that needs end-to-end knowledge cannot in general be implemented completely at a lower layer that lacks that knowledge~\citep{saltzerreedclark1984e2e}. Their setting was where to place \emph{functionality} in a distributed system; the setting here is where to place \emph{enforcement} over engineering work. The design instinct is the same --- put the deciding mechanism where the semantics needed to decide the property are available.
\end{quote}

\begin{figure}[htbp]
\centering
\includegraphics[width=\linewidth]{book/latex/assets/semantic-gap-levels-bf0468495b.pdf}
\caption{\emph{The semantic gap.} Different obligations become decidable at different boundaries. Model--implementation correspondence may be required at task completion even when intermediate commits legitimately disagree; rendered-page correctness requires an integrated execution that earlier source, build, and API checks cannot establish. Place each check at the earliest boundary where sufficient evidence is available.}
\label{semantic-gap}
\end{figure}

% <!-- point: model-drift-at-the-commit-is-the-wrong-granularity | Checking model drift at the commit is the wrong granularity; the property belongs to the completed unit of work. | terms: model-drift, semantic-gap -->

% <!-- index-def: semantic-gap -->

% <!-- index-example: model-drift -->

Consider an agent changing a system's call graph. The architectural model must agree with the implementation when the task is complete, but several intermediate commits may legitimately leave them inconsistent. Requiring correspondence after every commit would reject valid intermediate work. At task completion, the accumulated code changes and model updates can be checked together. The obligation belongs to the completed task, not to each commit.

% <!-- point: a-deployment-can-pass-every-check-and-still-serve-a-blank-page | Test the served page in a browser at the earliest boundary where the relevant composition can be exercised; a successful HTTP fetch alone does not establish correct rendering. | terms: semantic-gap, earliest-legible-boundary -->

A different obligation requires the assembled system to behave correctly. A deployment can pass source, build, and API checks yet serve a blank page. Those checks establish useful properties, but not that the application displays the expected content. A browser smoke test can load the page and check its rendered output. The test belongs at the earliest boundary where the integrated application can be exercised, such as a preview or staging deployment. A successful HTTP response alone does not establish that the page rendered correctly.

% <!-- gloss: Epic | A tracked unit of work larger than one commit: a feature or effort with its own file, plan, and definition of done, carried out over several agent dispatches. -->

% <!-- gloss: Definition of done (DoD) | The explicit list of conditions a unit of work must meet to count as finished — tests present, docs updated, the model in sync — checked at the end rather than assumed. -->

\section{The Earliest Decidable Boundary}\label{heading:the-earliest-decidable-boundary}

% <!-- section-terms: earliest-legible-boundary, semantic-gap -->

% <!-- point: place-enforcement-at-the-earliest-boundary-where-the-obligation-is-legible-and-decidable | The design rule: place enforcement at the earliest boundary where the obligation is both legible and decidable. | terms: earliest-legible-boundary, semantic-gap -->

% <!-- index-def: earliest-legible-boundary -->

Both examples illustrate the same rule for placing enforcement:

\begin{quote}
\textbf{Place enforcement at the earliest boundary where the obligation is both legible and decidable.}
\end{quote}

\emph{Earliest} matters because delay makes a failure more expensive to repair. \emph{Legible} matters because moving the check earlier than the property permits produces false refusals or meaningless checks. The rule explains the whole mechanism menu at a glance:

\begin{itemize}
\tightlist
\item
  \textbf{Static type checks} evaluate expressions against declared types before execution, once the relevant code and type information are available.
\item
  \textbf{Action hooks} act at one tool call, one command --- where a single move is legible.
\item
  \textbf{Task-completion checks} act at the work-unit boundary --- where "done" first becomes readable across a body of commits.
\item
  \textbf{Deployment gates} act after broader evidence has accumulated --- the last boundary before the artifact is accepted.
\end{itemize}

And some properties stay illegible until later still: a leak visible only across ten calls, a cost overrun visible only across a season of jobs. No boundary makes them readable sooner.

% <!-- point: the-part-two-rhyme-coarsest-representation-earliest-boundary | The rhyme: Modeling chooses the reduction that answers the question without carrying unnecessary detail; Alignment chooses the earliest boundary where the obligation is legible. | terms: semantic-gap, earliest-legible-boundary, modeling-principle, alignment-principle -->

\magevalue{chapter:modeling} and \magevalue{chapter:alignment} therefore use companion selection rules:

% <!-- inline-quote -->

\begin{quote}
\textbf{Modeling} --- \emph{choose the reduction that makes the engineering \textbf{question} answerable without carrying unnecessary detail.}

\textbf{Alignment} --- \emph{choose the earliest boundary at which the \textbf{obligation} becomes legible and decidable.}
\end{quote}

Both rules keep irrelevant detail from driving the engineering decision.

\begin{itemize}
\tightlist
\item
  Modeling avoids carrying detail the question does not need.
\item
  Alignment avoids enforcing an obligation before the necessary evidence exists.
\end{itemize}

Choose the wrong abstraction in \magevalue{chapter:modeling} and the required property is not represented. Choose the wrong boundary here and the required evidence is not available.

% <!-- point: real-boundaries-are-sometimes-approximate-a-roughly-right-hook-beats-the-perfect-one | Real boundaries are sometimes approximate: a hook that fires early enough to leave budget for the work beats the theoretically ideal hook. | terms: earliest-legible-boundary, semantic-gap -->

Real boundaries sometimes require a margin for the mechanism to act. Suppose an agent must preserve a durable hand-off before its working context is compacted. Waiting until no room remains may leave too little capacity to write it. A practical trigger can act earlier, when context is becoming full and the durable record is stale. The trigger approximates the desired boundary while preserving enough time and capacity to complete the required action.

% <!-- point: the-chapter-supplied-both-requirements-next-question-is-what-the-mechanism-does | The section has supplied both requirements for enforcement; the next question is what the mechanism at that boundary actually does. | terms: independent-obligation, semantic-gap -->

With an independent obligation and a boundary where it becomes legible, the next question is what the mechanism at that boundary actually does.
